%PDFTexify
\RequirePackage{atbegshi}
%\documentclass[envcountsect]{beamer}
\documentclass[pdf, 10pt]{beamer}
\usepackage{tikz}
\usepackage{beamerthemesplit}
\usepackage[german]{babel} 
\usepackage[utf8]{inputenc} %ACHTUNG: ALLE .tex-Dateien müssen UTF8-Codiert sein!
\usepackage{pgf,pgfarrows,pgfnodes,pgfautomata,pgfheaps}
\usepackage{subfigure,graphicx}
\usepackage{amsmath,amssymb}
\usepackage{color}
\usepackage{enumerate}
\usepackage{rotating}
\usepackage{pgfplots}
\usepackage{textpos}
\usepackage{xspace}
\usepackage{dsfont}
\usepackage{ifthen}
\usepackage{array}
\usepackage{multirow}
\usepackage{eurosym}
\usepackage{etoolbox}
\usepackage{calc}
\usepackage[export]{adjustbox}
\usepackage{xargs}
\usepackage{fancyvrb}
\usepackage[linesnumbered,ruled]{algorithm2e}
\usepackage{paralist} %für compactenum
\newboolean{printlogo}

%widerspruchsblitz
\usepackage{stmaryrd}

%tikz
\usepackage{pgf}
\usepackage{graphics}
%\usepackage[margin=1in]{geometry}
\usetikzlibrary{calc,arrows,automata}
 %to enable backgrounds in tikz graphics -> edges do not overlay vertices
\usetikzlibrary{backgrounds}
\pgfdeclarelayer{myback}
\pgfsetlayers{background,myback,main}



%\mode<presentation>{\usetheme{Warsaw}}
%\setbeamertemplate{navigation symbols}{}%remove navigation symbols
\beamertemplatenavigationsymbolsempty



%----------------------------------------------------------------------------
\newcommand{\vortragstitel}{Der Nullstellensatz in der Diskreten Mathematik} %Hier den Vortragstitel einfügen
\newcommand{\currentsemester}{SS 2016}  %aktuelles Semester
\newcommand{\lecturetitle}{Proseminar Diskrete Mathematik} % LV-Titel
\newcommand{\lecturesubtitle}{}
\newcommand{\lecturer}{Anna Jellen, Kerstin Maier} % Vortragende
\newcommand{\email}{kerstinma@edu.aau.at} % email-Kontakt
\newcommand{\tocTitle}{Inhalt} % Beschriftung (Titel) der Inhaltsfolien
\newcommand{\refTitle}{Literatur} % Beschriftung Literaturfolien
\setboolean{printlogo}{true}    % Logo auf den Folien ein/ausschalten
%----------------------------------------------------------------------------

%----------------------------------------------------------------------------
% Farben für die Folien (blaue und rote) in RGB-Notation
\definecolor{UniBlue}{RGB}{0,0,128}
\definecolor{UniRed}{RGB}{128,0,0}
%----------------------------------------------------------------------------

%input "literatur.bib"
\setbeamertemplate{bibliography item}[text]

\newcommand{\uebung}{\blau{Übung}}
\newcommand{\rb}[1]{\raisebox{1.5ex}[-1.5ex]{#1}}
\newcommand{\MOD}{\mathop{\text{MOD}}}
\newcommand{\TO}{\Rightarrow}
\newcommand{\DIV}{\mathop{\text{DIV}}}
%\newcommand{\ggt}{\text{ggT}}
\newcommand{\dividesnot}{\!\!\not|}
\newcommand{\R}{\mathds{R}}
\newcommand{\Z}{\mathds{Z}}
\newcommand{\N}{\mathds{N}}
\newcommand{\To}{\rightarrow}
\newcommand{\eps}{\varepsilon}
\newcommand{\set}[1]{\left\{#1\right\}}
\newcommand{\abs}[1]{\left|#1\right|}
\newcommand{\norm}[1]{\left\|#1\right\|}
\newcommand{\I}{\mathcal{I}}
\newcommand{\compl}[1]{\overline{#1}}

\newcommand{\rot}[1]{\textcolor[rgb]{1.00,0.00,0.00}{#1}\xspace}
\newcommand{\blau}[1]{\textcolor[rgb]{0.00,0.00,1.00}{#1}\xspace}
\newcommand{\gruen}[1]{\textcolor[rgb]{0.00,1.00,0.00}{#1}\xspace}
\newcommand{\grau}[1]{\textcolor[rgb]{0.50,0.50,0.50}{#1}\xspace}
\newcommand{\blueColor}{
    \setbeamercolor{title}{fg=UniBlue}
    \setbeamercolor{frametitle}{fg=UniBlue}
    \setbeamercolor{structure}{fg=UniBlue}
}
\newcommand{\redColor}{
    \setbeamercolor{title}{fg=UniRed}
    \setbeamercolor{frametitle}{fg=UniRed}
    \setbeamercolor{structure}{fg=UniRed}
}

\newcommand\blfootnote[1]{%
  \begingroup
  \renewcommand\thefootnote{}\footnote{#1}%
  \addtocounter{footnote}{-1}%
  \endgroup
}


\newboolean{onlyonerefslide}
\newcommand{\singlerefslide}{\setboolean{onlyonerefslide}{true}}
\newenvironment{lect}[2]
{
\setcounter{framenumber}{0}
\addtocounter{section}{1}
\part{#1}
\lecture{#1}{#2}
\addtocounter{section}{-1}
\section{#1}\label{#2}
\begin{slide}{\tocTitle}\tableofcontents[currentsubsection,subsectionstyle=show]\end{slide}
\setboolean{onlyonerefslide}{false}
}
%{
%\bibliographystyle{alpha}
%\ifthenelse{\boolean{onlyonerefslide}}{\begin{slide}{\refTitle}
%\bibliography{literatur}
%\end{slide}}{\begin{mslide}{\refTitle}
%\bibliography{literatur}
%\end{mslide}}
%}

\newcommand{\emptypage}{\frame{}}
\newtheorem{satz}{Satz}[section]
\newtheorem{thm}{Theorem}
\newtheorem{kor}{Korollar}
\newtheorem{ex}{Beispiel}
\newtheorem{lem}{Lemma}

\newcounter{bsp}
\newenvironment{bsp}[1][]{\refstepcounter{bsp}\par\medskip\noindent%
   \blau{Beispiel~\thebsp #1:}}{\hfill$\diamond$\medskip}
\numberwithin{bsp}{section}

\renewcommand{\qed}{\hfill$\square$}
\newcommand{\bew}{\underline{Beweis}:~}
\newenvironmentx{beweis}[1][1=\empty]{\ifthenelse{\equal{#1}{\empty}}{\par\medskip\noindent\bew}
{\par\medskip\noindent\underline{Beweis} (#1):~}}{\qed}

\newenvironment{slide}[1]{\blueColor\begin{frame}\frametitle{#1}}{\end{frame}}
\newenvironment{mslide}[1]{\blueColor\begin{frame}[allowframebreaks]\frametitle{#1}}{\end{frame}}
\newenvironment{redslide}[1]{\redColor\begin{frame}\frametitle{#1}}{\end{frame}}
\newenvironment{redmslide}[1]{\redColor\begin{frame}[allowframebreaks]\frametitle{#1}}{\end{frame}}

\usetheme{Boadilla}
\usecolortheme{seahorse}
\setbeamertemplate{blocks}[rounded][shadow=true]
\setbeamertemplate{caption}[numbered]


\setbeamertemplate{footline}
{
  \leavevmode%
  \hbox{%
  \begin{beamercolorbox}[wd=.7\paperwidth,ht=2.25ex,dp=1ex,left]{author in head/foot}%
    \usebeamerfont{author in head/foot}\space\hspace{1cm}\lecturer\hspace{1cm}\lecturetitle%\insertshortauthor
  \end{beamercolorbox}%
  \begin{beamercolorbox}[wd=.3\paperwidth,ht=2.25ex,dp=1ex,right]{title in head/foot}%
    \usebeamerfont{title in head/foot}\currentsemester%\insertshorttitle
    \hspace*{3em}
    \insertframenumber{}%/ \inserttotalframenumber
  \end{beamercolorbox}}%
  \vskip0pt%
}


\ifthenelse{\boolean{printlogo}}{
\addtobeamertemplate{frametitle}{}{%
\begin{textblock*}{100mm}(.85\textwidth,-0.75cm)
\includegraphics[width=2cm]{aaulogo}
\end{textblock*}
}
}{}

\AtBeginLecture{\author{\thesection\quad\insertlecture}\frame{\titlepage}}

\setbeamertemplate{theorems}[numbered]
\setbeamertemplate{caption}[numbered]
\numberwithin{figure}{section}

\title[\lecturetitle]{\lecturetitle}
\ifthenelse{\equal{\lecturesubtitle}{}}{
\author{\phantom{Yg}\lecturer\phantom{Yg}}
\institute{Institut für Mathematik\\Alpen-Adria-Universität Klagenfurt}
%\\\texttt{\email}}
\date{\currentsemester}
}
{%non-empty subtitle
\author{\phantom{Yg}\lecturesubtitle\phantom{Yg}}
\institute{\lecturer\\Institut für Mathematik \\Alpen-Adria-Universität Klagenfurt}
%\\\texttt{\email}}
\date{\currentsemester}
}

\begin{document}

\blueColor

%\frame{\titlepage}

\begin{lect}{\vortragstitel}{lec:vortrag}
\subsection{Einleitung}
\begin{mslide}{Einleitung}

\begin{definition}
Die \textcolor{blue}{k-Färbung} eines Graphen ist eine Abbildung $f: V(G) \rightarrow S$ mit $xy \in E(G): f(x) \neq f(y)$, wobei S Menge der Farben und $\vert S \vert = k$. 
\end{definition}

\begin{figure}[htp]
\begin{center}
\resizebox{.4\linewidth}{!}{
\begin{tikzpicture}[->,>=stealth',shorten >=1pt,auto,node distance=2.5cm,
                    semithick]
  \tikzstyle{every state}=[draw=black,text=black]
	\tikzstyle{every edge} = [draw,thick,-]

  \node[state,fill=red] 	(0)              	{1};
  \node[state,fill=blue]  (1) [above right of=0] 	{2};
  \node[state,fill=yellow] 	(2) [right of=0]  {3};
  \node[state,fill=blue]  (3) [below of=0] 	{4};
  \node[state,fill=yellow] 	(4) [right of=3]  {5};
  \node[state,fill=red]  (5) [right of=2] 	{6};
 

  
  \path[draw=black,thick] 
  		%% drone 1
  		(0)  edge node {} (2)
  		(3)  edge  node {} (4)
  		(5)  edge  node {} (4)
  		(1)  edge  node {} (0)
  		(2)  edge  node {} (3)
  		(1)  edge  node {} (5)
			(2)  edge  node {} (5)
			(3)  edge  node {} (5)
  		;
\end{tikzpicture}
}
\caption[3-Färbung]{3-Färbung}
  \label{Faerbung}
\end{center}
\end{figure}

\framebreak 

Sei $G = (V, E)$ ein Graph
\begin{align*}
&x_i^k - 1 = 0 \quad \forall i \in V, \\
& \sum_{s = 0}^{k-1} x_i^{k-1-s}x_j^s = 0 \quad \forall (i,j) \in E.
\end{align*} 

\begin{thm}
Graph G ist k-färbbar $\Leftrightarrow$ System besitzt eine komplexe Lösung
\end{thm}
\begin{thm}
G ist k-färbbar und k ungerade $\Leftrightarrow$ System besitzt gemeinsame Nullstelle über $\overline{\mathbb{F}_2}$, wobei  $\overline{\mathbb{F}_2}$ algebraische Abschluss über dem endlichen Körper mit zwei Elementen
\end{thm}
\nocite{*}

\framebreak

\begin{proof}
Angenommen die Aussage ist wahr über den komplexen Zahlen $\mathbb{C}$. \\
 "`$\Rightarrow$"' Sei G k-färbbar, ordne jeder Farbe die k-te Einheitswurzel zu. Sei die j-te Farbe $\beta_j = e^{2\Pi j/k}$; substituiere alle $x_l$ mit der zugehörigen Einheitswurzel der Farbe des l-ten Knotens. Also haben wir eine Lösung des Systems: Die Gleichungen $x_i^k-1 = 0$ sind trivialerweise erfüllt. \\Wir betrachten nun die Kantengleichungen: Wir nehmen eine Kante $ij$, da $x_i$ und $x_j$ durch Einheitswurzeln substituiert wurden, gilt $x_i^k - x_j^k = 0$. Des Weiteren gilt: 
\begin{align*}
x_i^k-x_j^k = (x_i-x_j)(x_i^{k-1}+x_i^{k-2}x_j+x_i^{k-3}x_j^2+\ldots+x_j^{k-1}) = 0;
\end{align*}   
Durch die Substitution mit unterschiedlichen Einheitswurzeln gilt $x_i - x_j \neq 0$, also muss der andere Faktor, der den Kantengleichungen entspricht, 0 sein.
\end{proof}

\framebreak 

\begin{proof}
"`$\Leftarrow$"' Angenommen die Gleichungen seien erfüllt, d.h. der Lösungspunkt muss aus k-ten Einheitswurzeln bestehen. Den benachbarten Knoten müssen verschiedene Einheitswurzeln zugeordnet werden, da: \\
Angenommen einem Paar benachbarter Knoten $ij$ wird die selbe Einheitswurzel zugewiesen. Die Gleichung $x_i^{k-1}+x_i^{k-2}x_j+x_i^{k-3}x_j^2+\ldots+x_j^{k-1} = 0$ wird dann zu $\beta^{k-1}+\beta^{k-1}+\ldots+\beta^{k-1} = k\beta^{k-1} = 0$, jedoch $\beta \neq 0 \lightning$ 
\end{proof}

\end{mslide}

\subsection{Hilberts Nullstellensatz}
\begin{mslide}{Hilberts Nullstellensatz}

\begin{block}{Problemdarstellung}
Gegeben: $f_1,\ldots,f_m \in \mathbb{K}[x_1,\ldots,x_n]$ \\
Gesucht: Lösung $x$ für das System $f_1(x) = 0, f_2(x) = 0, \ldots f_m(x) = 0$ (wird auch geschrieben als $F(x) = 0$) \\
Ziel ist es eine Lösung zu diesem System zu finden beziehungsweise zu zeigen, dass es keine Lösung gibt. \\
\end{block}


\begin{theorem}[Fredholm's Alternativtheorem] 
Das Lineare Gleichungssystem $Ax=b$ besitzt genau dann eine Lösung, wenn ein Vektor $y$ mit der Eigenschaft $y^TA=0^T$ und $y^Tb\neq 0^T$ existiert. 
\end{theorem}

\framebreak

\begin{theorem}[Hilbert's Nullstellensatz] 
Sei F = $\{f_1,\ldots,f_m\} \subseteq \mathbb{K}[x_1,\ldots,x_n]$.\\
Die Varietät $\{x \in \overline{\mathbb{K}^n} : f_1(x)=0,\ldots,f_s(x)=0\}$ ist genau dann leer, wenn 1 zum Ideal $\left\langle F \right\rangle$, das aus F generiert wurde, gehört. Man beachte $1 \in \left\langle F \right\rangle$ bedeutet, dass Polynome $\beta_1,\ldots,\beta_m$ im Ring $\mathbb{K}[x_1,\ldots,x_n]$ existieren, sodass $1 = \sum_{i=1}^m \beta_i f_i$. Diese polynomielle Identität ist ein Zertifikat für die Unlösbarkeit von $F(x) = 0$. 
\end{theorem}


\end{mslide}

%Falls mehr als eine Literaturfolie benötigt wird, bitte die folgende Anweisung entfernen oder auskommentieren
%\singlerefslide
\end{lect}

\end{document}
